% --- Archon physics formalization source begin ---
% archon:physics
% archon:covers PhyXMiniProblems/problem_phyx_mini_0522.lean
% archon:source-report reports/phyx_mini/problem_phyx_mini_0522.source.json

\chapter{Physics problem phyx\_mini\_0522}
\label{ch:PhyXMiniProblems_problem_phyx_mini_0522}

\paragraph{Problem source.}
\#\# Physical scenario

Van der Waals bonds arise from the interaction between two permanent or induced electric dipole moments in a pair of atoms or molecules. Consider two identical dipoles, each consisting of charges \textbackslash{}(+q\textbackslash{}) and \textbackslash{}(-q\textbackslash{}) separated by a distance \textbackslash{}(d\textbackslash{}) and oriented as shown in figure.

\#\# Question

Calculate the electric potential energy, expressed in terms of the electric dipole moment \textbackslash{}(p = qd\textbackslash{}), for the situation where \textbackslash{}(r \textbackslash{}gg d\textbackslash{}).

\#\# Answer choices

- A: A: \textbackslash{}( U = \textbackslash{}frac\{-2p\textasciicircum{}2\}\{3\textbackslash{}pi\textbackslash{}epsilon\_0 r\textasciicircum{}3\} - \textbackslash{}frac\{2p\textasciicircum{}2\}\{3\textbackslash{}pi\textbackslash{}epsilon\_0 d\textasciicircum{}3\} \textbackslash{})

- B: B: \textbackslash{}( U = \textbackslash{}frac\{-2p\textasciicircum{}2\}\{3\textbackslash{}pi\textbackslash{}epsilon\_0 r\textasciicircum{}3\} + \textbackslash{}frac\{2p\textasciicircum{}2\}\{3\textbackslash{}pi\textbackslash{}epsilon\_0 d\textasciicircum{}3\} \textbackslash{})

- C: C: \textbackslash{}( U = \textbackslash{}frac\{-2p\textasciicircum{}2\}\{4\textbackslash{}pi\textbackslash{}epsilon\_0 r\textasciicircum{}3\} + \textbackslash{}frac\{2p\textasciicircum{}2\}\{4\textbackslash{}pi\textbackslash{}epsilon\_0 d\textasciicircum{}3\} \textbackslash{})

- D: D: \textbackslash{}( U = \textbackslash{}frac\{-2p\textasciicircum{}2\}\{4\textbackslash{}pi\textbackslash{}epsilon\_0 r\textasciicircum{}3\} - \textbackslash{}frac\{2p\textasciicircum{}2\}\{4\textbackslash{}pi\textbackslash{}epsilon\_0 d\textasciicircum{}3\} \textbackslash{})

\#\# Figure caption (auxiliary; use the image as primary evidence)

The image depicts two sets of electric dipoles. Each dipole consists of a red sphere with a "+" sign labeled "+q" and a blue sphere with a "−" sign labeled "−q". The two spheres in each dipole are connected by a horizontal line. The dipoles are positioned horizontally, side by side, separated by a broken vertical line.

The distance between the center of the two dipoles is labeled "r" with an arrow indicating the measurement direction. The layout suggests a focus on the characteristics and interaction of these dipoles, often studied in physics and electromagnetism.

\#\# Dataset metadata

Quantum Phenomena; Physical Model Grounding Reasoning, Multi-Formula Reasoning

\paragraph{Recorded answer/context.}
D: D: \textbackslash{}( U = \textbackslash{}frac\{-2p\textasciicircum{}2\}\{4\textbackslash{}pi\textbackslash{}epsilon\_0 r\textasciicircum{}3\} - \textbackslash{}frac\{2p\textasciicircum{}2\}\{4\textbackslash{}pi\textbackslash{}epsilon\_0 d\textasciicircum{}3\} \textbackslash{})

\paragraph{Figure/image path.}
/root/proposal\_for\_physic/science-mango/phyx\_mini\_run/phyx\_data/test\_image/522.png

\paragraph{Formalization target.}
create a compiling Lean file with sorry bodies at `PhyXMiniProblems/problem\_phyx\_mini\_0522.lean`. The Lean declarations must preserve the physical quantities, dimensions or dimensional roles, figure labels, governing-law hypotheses, and final relation expressed by this problem.
Use Mathlib/Physlib names found through LeanExplore where available. If a physics API is missing, introduce faithful local abstractions rather than scalar placeholder aliases.

\begin{theorem}[Physics formalization target]
\label{thm:physics:phyx_mini_0522:target}
\lean{PhyXMiniProblems.ProblemPhyXMini0522.recorded_answer_choice_D}
\uses{def:physics:phyx-mini-0522:phyxminiproblems-problemphyxmini0522-collineardipolepairsetup, def:physics:phyx-mini-0522:phyxminiproblems-problemphyxmini0522-matchesproblemandprimaryfigure, def:physics:phyx-mini-0522:phyxminiproblems-problemphyxmini0522-hasphysicaldipoleparameters, def:physics:phyx-mini-0522:phyxminiproblems-problemphyxmini0522-isfarfieldregime, def:physics:phyx-mini-0522:phyxminiproblems-problemphyxmini0522-satisfiescoulombdipolelaws, def:physics:phyx-mini-0522:phyxminiproblems-problemphyxmini0522-iscorrectanswer}
The assigned autoformalize agent should translate this physics problem into Lean declarations in the covered file, with theorem and lemma proof bodies written as `by sorry`.
\end{theorem}
\begin{proof}
This is an autoformalization task, not a proof task. Produce statements that can later be proved by the physics prover without weakening the physical model.
\end{proof}

% --- Archon declaration topology begin ---
\section{Declaration topology}
\begin{definition}[Length Quantity]
  \label{def:physics:phyx-mini-0522:phyxminiproblems-problemphyxmini0522-lengthquantity}
  \lean{PhyXMiniProblems.ProblemPhyXMini0522.LengthQuantity}
  A nonnegative physical length
\end{definition}
\begin{proof}
  This is the declared model definition of Length Quantity.
\end{proof}
\begin{definition}[Charge Magnitude Quantity]
  \label{def:physics:phyx-mini-0522:phyxminiproblems-problemphyxmini0522-chargemagnitudequantity}
  \lean{PhyXMiniProblems.ProblemPhyXMini0522.ChargeMagnitudeQuantity}
  A nonnegative electric-charge magnitude
\end{definition}
\begin{proof}
  This is the declared model definition of Charge Magnitude Quantity.
\end{proof}
\begin{definition}[electric Dipole Moment Dimension]
  \label{def:physics:phyx-mini-0522:phyxminiproblems-problemphyxmini0522-electricdipolemomentdimension}
  \lean{PhyXMiniProblems.ProblemPhyXMini0522.electricDipoleMomentDimension}
  The physical dimension of an electric dipole moment, charge times length
\end{definition}
\begin{proof}
  This is the declared model definition of electric Dipole Moment Dimension.
\end{proof}
\begin{definition}[Electric Dipole Moment Magnitude Quantity]
  \label{def:physics:phyx-mini-0522:phyxminiproblems-problemphyxmini0522-electricdipolemomentmagnitudequantity}
  \lean{PhyXMiniProblems.ProblemPhyXMini0522.ElectricDipoleMomentMagnitudeQuantity}
  \uses{def:physics:phyx-mini-0522:phyxminiproblems-problemphyxmini0522-electricdipolemomentdimension}
  A nonnegative electric-dipole-moment magnitude
\end{definition}
\begin{proof}
  This is the declared model definition of Electric Dipole Moment Magnitude Quantity.
\end{proof}
\begin{definition}[vacuum Permittivity Dimension]
  \label{def:physics:phyx-mini-0522:phyxminiproblems-problemphyxmini0522-vacuumpermittivitydimension}
  \lean{PhyXMiniProblems.ProblemPhyXMini0522.vacuumPermittivityDimension}
  Vacuum permittivity has dimension C² T² M⁻¹ L⁻³, equivalently charge squared divided by energy times length
\end{definition}
\begin{proof}
  This is the declared model definition of vacuum Permittivity Dimension.
\end{proof}
\begin{definition}[Vacuum Permittivity Quantity]
  \label{def:physics:phyx-mini-0522:phyxminiproblems-problemphyxmini0522-vacuumpermittivityquantity}
  \lean{PhyXMiniProblems.ProblemPhyXMini0522.VacuumPermittivityQuantity}
  \uses{def:physics:phyx-mini-0522:phyxminiproblems-problemphyxmini0522-vacuumpermittivitydimension}
  A nonnegative physical vacuum permittivity
\end{definition}
\begin{proof}
  This is the declared model definition of Vacuum Permittivity Quantity.
\end{proof}
\begin{definition}[nonnegative SIReadout]
  \label{def:physics:phyx-mini-0522:phyxminiproblems-problemphyxmini0522-nonnegativesireadout}
  \lean{PhyXMiniProblems.ProblemPhyXMini0522.nonnegativeSIReadout}
  SI readout of a nonnegative dimensionful quantity
\end{definition}
\begin{proof}
  This is the declared model definition of nonnegative SIReadout.
\end{proof}
\begin{definition}[length In Meters]
  \label{def:physics:phyx-mini-0522:phyxminiproblems-problemphyxmini0522-lengthinmeters}
  \lean{PhyXMiniProblems.ProblemPhyXMini0522.lengthInMeters}
  \uses{def:physics:phyx-mini-0522:phyxminiproblems-problemphyxmini0522-lengthquantity, def:physics:phyx-mini-0522:phyxminiproblems-problemphyxmini0522-nonnegativesireadout}
  Length readout in metres
\end{definition}
\begin{proof}
  This is the declared model definition of length In Meters.
\end{proof}
\begin{definition}[charge Magnitude In Coulombs]
  \label{def:physics:phyx-mini-0522:phyxminiproblems-problemphyxmini0522-chargemagnitudeincoulombs}
  \lean{PhyXMiniProblems.ProblemPhyXMini0522.chargeMagnitudeInCoulombs}
  \uses{def:physics:phyx-mini-0522:phyxminiproblems-problemphyxmini0522-chargemagnitudequantity, def:physics:phyx-mini-0522:phyxminiproblems-problemphyxmini0522-nonnegativesireadout}
  Charge-magnitude readout in coulombs
\end{definition}
\begin{proof}
  This is the declared model definition of charge Magnitude In Coulombs.
\end{proof}
\begin{definition}[dipole Moment Magnitude In Coulomb Meters]
  \label{def:physics:phyx-mini-0522:phyxminiproblems-problemphyxmini0522-dipolemomentmagnitudeincoulombmeters}
  \lean{PhyXMiniProblems.ProblemPhyXMini0522.dipoleMomentMagnitudeInCoulombMeters}
  \uses{def:physics:phyx-mini-0522:phyxminiproblems-problemphyxmini0522-electricdipolemomentmagnitudequantity, def:physics:phyx-mini-0522:phyxminiproblems-problemphyxmini0522-nonnegativesireadout}
  Electric-dipole-moment readout in coulomb-metres
\end{definition}
\begin{proof}
  This is the declared model definition of dipole Moment Magnitude In Coulomb Meters.
\end{proof}
\begin{definition}[vacuum Permittivity In Farads Per Meter]
  \label{def:physics:phyx-mini-0522:phyxminiproblems-problemphyxmini0522-vacuumpermittivityinfaradspermeter}
  \lean{PhyXMiniProblems.ProblemPhyXMini0522.vacuumPermittivityInFaradsPerMeter}
  \uses{def:physics:phyx-mini-0522:phyxminiproblems-problemphyxmini0522-vacuumpermittivityquantity, def:physics:phyx-mini-0522:phyxminiproblems-problemphyxmini0522-nonnegativesireadout}
  Vacuum-permittivity readout in farads per metre
\end{definition}
\begin{proof}
  This is the declared model definition of vacuum Permittivity In Farads Per Meter.
\end{proof}
\begin{definition}[energy In Joules]
  \label{def:physics:phyx-mini-0522:phyxminiproblems-problemphyxmini0522-energyinjoules}
  \lean{PhyXMiniProblems.ProblemPhyXMini0522.energyInJoules}
  Signed energy readout in joules
\end{definition}
\begin{proof}
  This is the declared model definition of energy In Joules.
\end{proof}
\begin{definition}[Dipole Label]
  \label{def:physics:phyx-mini-0522:phyxminiproblems-problemphyxmini0522-dipolelabel}
  \lean{PhyXMiniProblems.ProblemPhyXMini0522.DipoleLabel}
  The two dipoles in their left-to-right order in the image
\end{definition}
\begin{proof}
  This is the declared model definition of Dipole Label.
\end{proof}
\begin{definition}[Dipole End]
  \label{def:physics:phyx-mini-0522:phyxminiproblems-problemphyxmini0522-dipoleend}
  \lean{PhyXMiniProblems.ProblemPhyXMini0522.DipoleEnd}
  The two endpoints of a horizontal dipole
\end{definition}
\begin{proof}
  This is the declared model definition of Dipole End.
\end{proof}
\begin{definition}[Charge Sign]
  \label{def:physics:phyx-mini-0522:phyxminiproblems-problemphyxmini0522-chargesign}
  \lean{PhyXMiniProblems.ProblemPhyXMini0522.ChargeSign}
  The sign carried by a point charge
\end{definition}
\begin{proof}
  This is the declared model definition of Charge Sign.
\end{proof}
\begin{definition}[Separation Label]
  \label{def:physics:phyx-mini-0522:phyxminiproblems-problemphyxmini0522-separationlabel}
  \lean{PhyXMiniProblems.ProblemPhyXMini0522.SeparationLabel}
  The distance symbols named by the prose and the primary image
\end{definition}
\begin{proof}
  This is the declared model definition of Separation Label.
\end{proof}
\begin{definition}[charge Sign Factor]
  \label{def:physics:phyx-mini-0522:phyxminiproblems-problemphyxmini0522-chargesignfactor}
  \lean{PhyXMiniProblems.ProblemPhyXMini0522.chargeSignFactor}
  \uses{def:physics:phyx-mini-0522:phyxminiproblems-problemphyxmini0522-chargesign}
  The scalar multiplier associated with a charge sign
\end{definition}
\begin{proof}
  This is the declared model definition of charge Sign Factor.
\end{proof}
\begin{definition}[Dipole Pair Figure]
  \label{def:physics:phyx-mini-0522:phyxminiproblems-problemphyxmini0522-dipolepairfigure}
  \lean{PhyXMiniProblems.ProblemPhyXMini0522.DipolePairFigure}
  \uses{def:physics:phyx-mini-0522:phyxminiproblems-problemphyxmini0522-dipolelabel, def:physics:phyx-mini-0522:phyxminiproblems-problemphyxmini0522-dipoleend, def:physics:phyx-mini-0522:phyxminiproblems-problemphyxmini0522-chargesign, def:physics:phyx-mini-0522:phyxminiproblems-problemphyxmini0522-separationlabel}
  Literal visual metadata from the primary raster. The r arrow is drawn between dashed center lines and contains a break indicating a long distance
\end{definition}
\begin{proof}
  This is the declared model definition of Dipole Pair Figure.
\end{proof}
\begin{definition}[Collinear Dipole Pair Setup]
  \label{def:physics:phyx-mini-0522:phyxminiproblems-problemphyxmini0522-collineardipolepairsetup}
  \lean{PhyXMiniProblems.ProblemPhyXMini0522.CollinearDipolePairSetup}
  \uses{def:physics:phyx-mini-0522:phyxminiproblems-problemphyxmini0522-lengthquantity, def:physics:phyx-mini-0522:phyxminiproblems-problemphyxmini0522-chargemagnitudequantity, def:physics:phyx-mini-0522:phyxminiproblems-problemphyxmini0522-electricdipolemomentmagnitudequantity, def:physics:phyx-mini-0522:phyxminiproblems-problemphyxmini0522-vacuumpermittivityquantity, def:physics:phyx-mini-0522:phyxminiproblems-problemphyxmini0522-dipolelabel, def:physics:phyx-mini-0522:phyxminiproblems-problemphyxmini0522-dipoleend, def:physics:phyx-mini-0522:phyxminiproblems-problemphyxmini0522-chargesign, def:physics:phyx-mini-0522:phyxminiproblems-problemphyxmini0522-dipolepairfigure}
  The physical system and its response quantities. The interaction profile is an SI scalar readout indexed by a center separation in metres. Its name records both the physical observable and its unit. The far-field coefficient likewise has the dimensional role joule-metres cubed. Neither field is assigned the coefficient sought in the question
\end{definition}
\begin{proof}
  This is the declared model definition of Collinear Dipole Pair Setup.
\end{proof}
\begin{definition}[dipole Center In Meters]
  \label{def:physics:phyx-mini-0522:phyxminiproblems-problemphyxmini0522-dipolecenterinmeters}
  \lean{PhyXMiniProblems.ProblemPhyXMini0522.dipoleCenterInMeters}
  \uses{def:physics:phyx-mini-0522:phyxminiproblems-problemphyxmini0522-dipolelabel}
  Center coordinate, in metres, when the left center is chosen as the origin
\end{definition}
\begin{proof}
  This is the declared model definition of dipole Center In Meters.
\end{proof}
\begin{definition}[endpoint Offset In Meters]
  \label{def:physics:phyx-mini-0522:phyxminiproblems-problemphyxmini0522-endpointoffsetinmeters}
  \lean{PhyXMiniProblems.ProblemPhyXMini0522.endpointOffsetInMeters}
  \uses{def:physics:phyx-mini-0522:phyxminiproblems-problemphyxmini0522-lengthinmeters, def:physics:phyx-mini-0522:phyxminiproblems-problemphyxmini0522-dipolelabel, def:physics:phyx-mini-0522:phyxminiproblems-problemphyxmini0522-dipoleend, def:physics:phyx-mini-0522:phyxminiproblems-problemphyxmini0522-collineardipolepairsetup}
  Signed endpoint offset from a dipole's center
\end{definition}
\begin{proof}
  This is the declared model definition of endpoint Offset In Meters.
\end{proof}
\begin{definition}[charge Position In Meters]
  \label{def:physics:phyx-mini-0522:phyxminiproblems-problemphyxmini0522-chargepositioninmeters}
  \lean{PhyXMiniProblems.ProblemPhyXMini0522.chargePositionInMeters}
  \uses{def:physics:phyx-mini-0522:phyxminiproblems-problemphyxmini0522-dipolelabel, def:physics:phyx-mini-0522:phyxminiproblems-problemphyxmini0522-dipoleend, def:physics:phyx-mini-0522:phyxminiproblems-problemphyxmini0522-collineardipolepairsetup, def:physics:phyx-mini-0522:phyxminiproblems-problemphyxmini0522-dipolecenterinmeters, def:physics:phyx-mini-0522:phyxminiproblems-problemphyxmini0522-endpointoffsetinmeters}
  Axial position of a labeled charge endpoint in the collinear model
\end{definition}
\begin{proof}
  This is the declared model definition of charge Position In Meters.
\end{proof}
\begin{definition}[point Charge Distance In Meters]
  \label{def:physics:phyx-mini-0522:phyxminiproblems-problemphyxmini0522-pointchargedistanceinmeters}
  \lean{PhyXMiniProblems.ProblemPhyXMini0522.pointChargeDistanceInMeters}
  \uses{def:physics:phyx-mini-0522:phyxminiproblems-problemphyxmini0522-dipolelabel, def:physics:phyx-mini-0522:phyxminiproblems-problemphyxmini0522-dipoleend, def:physics:phyx-mini-0522:phyxminiproblems-problemphyxmini0522-collineardipolepairsetup, def:physics:phyx-mini-0522:phyxminiproblems-problemphyxmini0522-chargepositioninmeters}
  Distance between two labeled point charges at a proposed center separation
\end{definition}
\begin{proof}
  This is the declared model definition of point Charge Distance In Meters.
\end{proof}
\begin{definition}[signed Charge In Coulombs]
  \label{def:physics:phyx-mini-0522:phyxminiproblems-problemphyxmini0522-signedchargeincoulombs}
  \lean{PhyXMiniProblems.ProblemPhyXMini0522.signedChargeInCoulombs}
  \uses{def:physics:phyx-mini-0522:phyxminiproblems-problemphyxmini0522-chargemagnitudeincoulombs, def:physics:phyx-mini-0522:phyxminiproblems-problemphyxmini0522-dipolelabel, def:physics:phyx-mini-0522:phyxminiproblems-problemphyxmini0522-dipoleend, def:physics:phyx-mini-0522:phyxminiproblems-problemphyxmini0522-chargesignfactor, def:physics:phyx-mini-0522:phyxminiproblems-problemphyxmini0522-collineardipolepairsetup}
  Signed SI charge at one endpoint of a dipole
\end{definition}
\begin{proof}
  This is the declared model definition of signed Charge In Coulombs.
\end{proof}
\begin{definition}[internal Coulomb Energy In Joules]
  \label{def:physics:phyx-mini-0522:phyxminiproblems-problemphyxmini0522-internalcoulombenergyinjoules}
  \lean{PhyXMiniProblems.ProblemPhyXMini0522.internalCoulombEnergyInJoules}
  \uses{def:physics:phyx-mini-0522:phyxminiproblems-problemphyxmini0522-dipolelabel, def:physics:phyx-mini-0522:phyxminiproblems-problemphyxmini0522-collineardipolepairsetup, def:physics:phyx-mini-0522:phyxminiproblems-problemphyxmini0522-pointchargedistanceinmeters, def:physics:phyx-mini-0522:phyxminiproblems-problemphyxmini0522-signedchargeincoulombs}
  The exact internal Coulomb energy of the two (+q,-q) pairs. Each unordered within-dipole pair occurs exactly once
\end{definition}
\begin{proof}
  This is the declared model definition of internal Coulomb Energy In Joules.
\end{proof}
\begin{definition}[cross Dipole Coulomb Energy In Joules]
  \label{def:physics:phyx-mini-0522:phyxminiproblems-problemphyxmini0522-crossdipolecoulombenergyinjoules}
  \lean{PhyXMiniProblems.ProblemPhyXMini0522.crossDipoleCoulombEnergyInJoules}
  \uses{def:physics:phyx-mini-0522:phyxminiproblems-problemphyxmini0522-dipoleend, def:physics:phyx-mini-0522:phyxminiproblems-problemphyxmini0522-collineardipolepairsetup, def:physics:phyx-mini-0522:phyxminiproblems-problemphyxmini0522-pointchargedistanceinmeters, def:physics:phyx-mini-0522:phyxminiproblems-problemphyxmini0522-signedchargeincoulombs}
  The exact cross-dipole Coulomb energy at center separation R. The nested sum contains the four pairs formed from one endpoint of each dipole
\end{definition}
\begin{proof}
  This is the declared model definition of cross Dipole Coulomb Energy In Joules.
\end{proof}
\begin{definition}[Matches Problem And Primary Figure]
  \label{def:physics:phyx-mini-0522:phyxminiproblems-problemphyxmini0522-matchesproblemandprimaryfigure}
  \lean{PhyXMiniProblems.ProblemPhyXMini0522.MatchesProblemAndPrimaryFigure}
  \uses{def:physics:phyx-mini-0522:phyxminiproblems-problemphyxmini0522-collineardipolepairsetup}
  Data read from the prose and primary image. In particular, both dipoles have +q on the left and -q on the right, and the arrow labeled r measures the center-to-center separation. This predicate contains no energy formula
\end{definition}
\begin{proof}
  This is the declared model definition of Matches Problem And Primary Figure.
\end{proof}
\begin{definition}[Has Physical Dipole Parameters]
  \label{def:physics:phyx-mini-0522:phyxminiproblems-problemphyxmini0522-hasphysicaldipoleparameters}
  \lean{PhyXMiniProblems.ProblemPhyXMini0522.HasPhysicalDipoleParameters}
  \uses{def:physics:phyx-mini-0522:phyxminiproblems-problemphyxmini0522-lengthinmeters, def:physics:phyx-mini-0522:phyxminiproblems-problemphyxmini0522-chargemagnitudeincoulombs, def:physics:phyx-mini-0522:phyxminiproblems-problemphyxmini0522-dipolemomentmagnitudeincoulombmeters, def:physics:phyx-mini-0522:phyxminiproblems-problemphyxmini0522-vacuumpermittivityinfaradspermeter, def:physics:phyx-mini-0522:phyxminiproblems-problemphyxmini0522-collineardipolepairsetup}
  Positivity and non-overlap conditions for the physical parameters
\end{definition}
\begin{proof}
  This is the declared model definition of Has Physical Dipole Parameters.
\end{proof}
\begin{definition}[Is Far Field Regime]
  \label{def:physics:phyx-mini-0522:phyxminiproblems-problemphyxmini0522-isfarfieldregime}
  \lean{PhyXMiniProblems.ProblemPhyXMini0522.IsFarFieldRegime}
  \uses{def:physics:phyx-mini-0522:phyxminiproblems-problemphyxmini0522-lengthinmeters, def:physics:phyx-mini-0522:phyxminiproblems-problemphyxmini0522-collineardipolepairsetup}
  A quantitative witness for the qualitative condition r ≫ d. The modeler chooses a dimensionless tolerance smaller than one, and each ratio d/r must lie below it. No numerical tolerance is invented from the raster
\end{definition}
\begin{proof}
  This is the declared model definition of Is Far Field Regime.
\end{proof}
\begin{definition}[Satisfies Coulomb Dipole Laws]
  \label{def:physics:phyx-mini-0522:phyxminiproblems-problemphyxmini0522-satisfiescoulombdipolelaws}
  \lean{PhyXMiniProblems.ProblemPhyXMini0522.SatisfiesCoulombDipoleLaws}
  \uses{def:physics:phyx-mini-0522:phyxminiproblems-problemphyxmini0522-lengthinmeters, def:physics:phyx-mini-0522:phyxminiproblems-problemphyxmini0522-chargemagnitudeincoulombs, def:physics:phyx-mini-0522:phyxminiproblems-problemphyxmini0522-dipolemomentmagnitudeincoulombmeters, def:physics:phyx-mini-0522:phyxminiproblems-problemphyxmini0522-vacuumpermittivityinfaradspermeter, def:physics:phyx-mini-0522:phyxminiproblems-problemphyxmini0522-energyinjoules, def:physics:phyx-mini-0522:phyxminiproblems-problemphyxmini0522-collineardipolepairsetup, def:physics:phyx-mini-0522:phyxminiproblems-problemphyxmini0522-internalcoulombenergyinjoules, def:physics:phyx-mini-0522:phyxminiproblems-problemphyxmini0522-crossdipolecoulombenergyinjoules}
  Coulomb's law, the definition p = qd, and the convention used to extract a far-field prediction from the exact interaction profile. The asymptotic field merely says that the stored coefficient is the limit of R³ U\_interaction(R). It does not state that limit's value. Similarly, the assembly field only says that the far-field prediction consists of the exact internal energy plus the extracted inverse-cube term
\end{definition}
\begin{proof}
  This is the declared model definition of Satisfies Coulomb Dipole Laws.
\end{proof}
\begin{definition}[Answer Choice]
  \label{def:physics:phyx-mini-0522:phyxminiproblems-problemphyxmini0522-answerchoice}
  \lean{PhyXMiniProblems.ProblemPhyXMini0522.AnswerChoice}
  Labels printed beside the four multiple-choice formulas
\end{definition}
\begin{proof}
  This is the declared model definition of Answer Choice.
\end{proof}
\begin{definition}[displayed Answer Energy In Joules]
  \label{def:physics:phyx-mini-0522:phyxminiproblems-problemphyxmini0522-displayedanswerenergyinjoules}
  \lean{PhyXMiniProblems.ProblemPhyXMini0522.displayedAnswerEnergyInJoules}
  \uses{def:physics:phyx-mini-0522:phyxminiproblems-problemphyxmini0522-lengthinmeters, def:physics:phyx-mini-0522:phyxminiproblems-problemphyxmini0522-dipolemomentmagnitudeincoulombmeters, def:physics:phyx-mini-0522:phyxminiproblems-problemphyxmini0522-vacuumpermittivityinfaradspermeter, def:physics:phyx-mini-0522:phyxminiproblems-problemphyxmini0522-collineardipolepairsetup, def:physics:phyx-mini-0522:phyxminiproblems-problemphyxmini0522-answerchoice}
  The energy expression printed for each answer choice, evaluated using the common dipole data represented by the left dipole
\end{definition}
\begin{proof}
  This is the declared model definition of displayed Answer Energy In Joules.
\end{proof}
\begin{definition}[Is Correct Answer]
  \label{def:physics:phyx-mini-0522:phyxminiproblems-problemphyxmini0522-iscorrectanswer}
  \lean{PhyXMiniProblems.ProblemPhyXMini0522.IsCorrectAnswer}
  \uses{def:physics:phyx-mini-0522:phyxminiproblems-problemphyxmini0522-energyinjoules, def:physics:phyx-mini-0522:phyxminiproblems-problemphyxmini0522-collineardipolepairsetup, def:physics:phyx-mini-0522:phyxminiproblems-problemphyxmini0522-answerchoice, def:physics:phyx-mini-0522:phyxminiproblems-problemphyxmini0522-displayedanswerenergyinjoules}
  A displayed choice is correct when it equals the modeled far-field energy
\end{definition}
\begin{proof}
  This is the declared model definition of Is Correct Answer.
\end{proof}
\begin{lemma}[internal Potential Energy formula]
  \label{lem:physics:phyx-mini-0522:phyxminiproblems-problemphyxmini0522-internalpotentialenergy-formula}
  \lean{PhyXMiniProblems.ProblemPhyXMini0522.internalPotentialEnergy_formula}
  \uses{def:physics:phyx-mini-0522:phyxminiproblems-problemphyxmini0522-lengthinmeters, def:physics:phyx-mini-0522:phyxminiproblems-problemphyxmini0522-dipolemomentmagnitudeincoulombmeters, def:physics:phyx-mini-0522:phyxminiproblems-problemphyxmini0522-vacuumpermittivityinfaradspermeter, def:physics:phyx-mini-0522:phyxminiproblems-problemphyxmini0522-energyinjoules, def:physics:phyx-mini-0522:phyxminiproblems-problemphyxmini0522-collineardipolepairsetup, def:physics:phyx-mini-0522:phyxminiproblems-problemphyxmini0522-matchesproblemandprimaryfigure, def:physics:phyx-mini-0522:phyxminiproblems-problemphyxmini0522-hasphysicaldipoleparameters, def:physics:phyx-mini-0522:phyxminiproblems-problemphyxmini0522-satisfiescoulombdipolelaws}
  The two attractive within-dipole Coulomb pairs give the constant internal contribution appearing in the recorded answer
\end{lemma}
\begin{proof}
  The conclusion follows from the governing relations and figure readouts named in the statement of internal Potential Energy formula.
\end{proof}
\begin{lemma}[far Field Interaction Coefficient formula]
  \label{lem:physics:phyx-mini-0522:phyxminiproblems-problemphyxmini0522-farfieldinteractioncoefficient-formula}
  \lean{PhyXMiniProblems.ProblemPhyXMini0522.farFieldInteractionCoefficient_formula}
  \uses{def:physics:phyx-mini-0522:phyxminiproblems-problemphyxmini0522-dipolemomentmagnitudeincoulombmeters, def:physics:phyx-mini-0522:phyxminiproblems-problemphyxmini0522-vacuumpermittivityinfaradspermeter, def:physics:phyx-mini-0522:phyxminiproblems-problemphyxmini0522-collineardipolepairsetup, def:physics:phyx-mini-0522:phyxminiproblems-problemphyxmini0522-matchesproblemandprimaryfigure, def:physics:phyx-mini-0522:phyxminiproblems-problemphyxmini0522-hasphysicaldipoleparameters, def:physics:phyx-mini-0522:phyxminiproblems-problemphyxmini0522-satisfiescoulombdipolelaws}
  The exact four-pair interaction profile has inverse-cube coefficient -2 p²/(4π ε₀) for the parallel collinear orientation in the image
\end{lemma}
\begin{proof}
  The conclusion follows from the governing relations and figure readouts named in the statement of far Field Interaction Coefficient formula.
\end{proof}
\begin{theorem}[collinear identical dipoles far field energy]
  \label{thm:physics:phyx-mini-0522:phyxminiproblems-problemphyxmini0522-collinear-identical-dipoles-far-field-energy}
  \lean{PhyXMiniProblems.ProblemPhyXMini0522.collinear_identical_dipoles_far_field_energy}
  \uses{def:physics:phyx-mini-0522:phyxminiproblems-problemphyxmini0522-lengthinmeters, def:physics:phyx-mini-0522:phyxminiproblems-problemphyxmini0522-dipolemomentmagnitudeincoulombmeters, def:physics:phyx-mini-0522:phyxminiproblems-problemphyxmini0522-vacuumpermittivityinfaradspermeter, def:physics:phyx-mini-0522:phyxminiproblems-problemphyxmini0522-energyinjoules, def:physics:phyx-mini-0522:phyxminiproblems-problemphyxmini0522-collineardipolepairsetup, def:physics:phyx-mini-0522:phyxminiproblems-problemphyxmini0522-matchesproblemandprimaryfigure, def:physics:phyx-mini-0522:phyxminiproblems-problemphyxmini0522-hasphysicaldipoleparameters, def:physics:phyx-mini-0522:phyxminiproblems-problemphyxmini0522-isfarfieldregime, def:physics:phyx-mini-0522:phyxminiproblems-problemphyxmini0522-satisfiescoulombdipolelaws}
  For r ≫ d, the far-field electric potential energy is the attractive inverse-cube dipole interaction plus the two internal binding energies. This is the formula printed as answer D. Blueprint
\end{theorem}
\begin{proof}
  The conclusion follows from the governing relations and figure readouts named in the statement of collinear identical dipoles far field energy.
\end{proof}
% --- Archon declaration topology end ---
% --- Archon physics formalization source end ---
